\begin{abstract}
\noindent
A survey covering many regions produces one estimate for each. Those
estimates are then used to describe a place the survey did not cover: a region
newly added to a frame, or one too small to sample. That means separating real
variation between places from sampling error, and the sampling error is itself
estimated, from designs supplying few degrees of freedom. We ask what that
costs.

The answer is not about estimation accuracy. Competing methods can be given
exactly the same information about the unknown spread of the true values and
still report very different widths, because they differ in how strongly the
reported width reacts to being wrong about that spread. We define that reaction,
show it is weaker for a method that subtracts the sampling error than for methods
whose width is proportional to the square root of a variance bound, and prove an
exact identity relating the two. Three things follow. The number of places a
practitioner needs is far smaller than a requirement stated on the variance alone
implies, and this is a property of the problem, not of the estimator. Two methods can change places, one narrower in one survey
and wider in another, at a computable point. And in one matched case the width
ratio depends on nothing but the number of places and the error budget; we
recorded that prediction before testing it, and one of its three parts failed.

We then give an interval exploiting the weaker reaction, with a guarantee stated
under a normal model and a bound computed rather than searched for, and report
where it stops working. On three European Social Survey questions
across three rounds it is narrower than a comparable method in 44 of 45 configurations,
under stated approximations and with no claim about coverage.
\end{abstract}
